To isolate final evidence selection, we replayed the frozen GPT-4o trajectories for 80 injected variants from 20 source SCCs. The control applies LEA's existing aggregation score to the same recorded local outputs and selects exactly as many nodes as the Full core in each case. Full and the LEA-score control achieve Risk-all of 75.00\% and 72.50\%, respectively (Full minus control: +2.50 percentage points; 95\% paired source-SCC bootstrap interval [-16.25, 20.00]). Both retain 8.10 nodes on average, with 51.26\% compression. The paired interval does not establish an advantage for either selector on these fixed trajectories. The comparison conditions on evidence acquired by Full and does not isolate search-policy effects or online replacement benefits; it requires no new model calls.

The prespecified mean-risk top-$k$ diagnostic achieved 93.75\% Risk-all at the same per-case output size, compared with 75.00\% for Full (diagnostic minus Full: $+18.75$ percentage points; 95\% paired source-SCC interval $[7.50,32.50]$). This diagnostic was fixed before computing the new selector outcomes. It shows that the current composite core selector can omit endpoints already recoverable from the acquired node scores; the results therefore do not support an independent advantage of that selector on injected-endpoint retention. This offline finding does not establish improved expert-evidence quality, conflict-absence detection, or a superior independently deployed acquisition policy.
